% !TeX root = ../main.tex
% ============================================================
%  sections/06_computational_experiments.tex
% ============================================================

\section{Computational Experiments}
\label{sec:experiments}

We use the model and the \DBP of \Cref{sec:model,sec:methodology} to answer
four questions. \emph{What} carbon tax does the regulator choose, and what
cost--emission trade-off does it induce (\Cref{sec:exp:tradeoff})?
\emph{Which} model inputs drive cost and emissions
(\Cref{sec:exp:sensitivity})? \emph{Which} design levers --- cleaner
technology or transport --- actually abate emissions
(\Cref{sec:exp:scenarios})? And does the \DBP deliver these answers reliably
(\Cref{sec:exp:scalability})? All reported numbers are regenerated from the
saved solver outputs by \path{tools.experiments.sensitivity analyze-all} and
stored under \path{results/analysis}; the follower-objective identity
$\Cost+\tau\,\Emis=$ follower objective holds across all runs (maximum
residual $6\times10^{-8}$).

\subsection{Experimental setup}
\label{sec:exp:setup}

\paragraph{Instances.}
The benchmark is produced by the repository generator. Twelve base codes are
defined as the Cartesian product of two network sizes (I17, I25), two
material/product scales (M4, M8), and three planning horizons (T3, T6, T12),
under the naming convention \texttt{STD\_I\textit{n}\_M\textit{m}\_T\textit{t}}
(\Cref{tab:instances}). The eight T3 and T6 instances are solved to
validation and form the basis of the analysis; the four T12 instances are
intractable within our single-machine compute budget and are excluded
throughout. The leader's tax grid spans $\tau\in[0,200]$ USD/t\,CO$_2$e
(index $\taxindex\in\{0,\ldots,20\}$, \Cref{sec:model:variables}); this range
is justified against current and prospective real-world carbon prices in
\Cref{app:data:tax_range}.

\begin{table}[ht]
  \centering
  \caption{Generated benchmark-instance codes.}
  \label{tab:instances}
  % !TeX root = ../main.tex
% ============================================================
%  tables/tab_instances.tex
% ============================================================

\renewcommand{\arraystretch}{1.2}
\begin{tabular}{@{}llp{0.52\linewidth}@{}}
  \toprule
  \textbf{Code} & \textbf{Values} & \textbf{Implemented meaning} \\
  \midrule
  $I$ & 17, 25
    & Total entity nodes in
      $\setS\cup\setF\cup\setW\cup\setR\cup\mathrm{AIR}\cup\mathrm{PORT}$.
      I17 uses 3 suppliers, 3 factories, 3 warehouses, 4 retailers,
      2 airports, and 2 seaports; I25 increases warehouses to 9 and
      airports to 4. \\
  $M$ & 4, 8
    & Material/product scale. M4 uses 2 raw materials and 1 final
      product; M8 uses 4 raw materials and 2 final products. \\
  $T$ & 3, 6, 12
    & Number of planning periods in the generated instance. \\
  Modes & \makecell[l]{mrd1--mrd2 ($\setMrd$),\\ma1--ma2 ($\setMa$),\\ms1--ms2 ($\setMs$)}
    & Common transport catalogue with two alternatives per family,
      following Tautenhain et al. The road modes ($\setMrd$) are
      company-owned truck alternatives; the air ($\setMa$) and sea
      ($\setMs$) modes are outsourced carrier alternatives. \\
  $\taxindex$ & 0,\ldots,20
    & Discrete tax index, used only in the leader's formulation and the
      \DBP (\Cref{sec:model:variables}). The economically meaningful
      rate is $\tau=10\,\taxindex$ USD/t\,CO\textsubscript{2}e, so the
      implemented domain is $[0,200]$ USD/t (justified in
      \Cref{app:data:tax_range}); all reported results use $\tau$. \\
  \bottomrule
\end{tabular}

\end{table}

\paragraph{Scenario design.}
To probe the supply chain's design levers, each T3 base instance is expanded in three
ways: a \emph{technology} scenario that enlarges the production and
remanufacturing technology menu, a \emph{transport} scenario that enlarges
the transport-mode catalogue, and a \emph{combined} scenario that applies
both. This yields a $2\times2$ design (technology on/off $\times$ transport
on/off) per instance.

\paragraph{Sensitivity design.}
A one-at-a-time (OAT) screen perturbs eleven scalar input factors --- together
with several structural variants (\eg removal of the mandatory-return floor)
--- around four base instances, for 75 sensitivity runs in total. Each output
change is normalized to an elasticity so that factors are comparable. The
screen is complemented by a Morris elementary-effects design
\citep{morris1991,saltelli2010} on \texttt{STD\_I17\_M4\_T3} and by extended
tax-range runs that lift the leader
ceiling from 200 USD/t to 400 and 800 USD/t, used in \Cref{sec:exp:tradeoff} to
test whether the base range truncates the regulator's optimum.

\paragraph{Solver configuration.}
All runs use the lexicographic exact workflow \texttt{lf\_lex\_exact}
described in \Cref{sec:method:implementation}, with relative \DBP tolerance
$10^{-4}$, per-subproblem Gurobi MIP gap $10^{-4}$, six threads, and no time
limit; base runs are cross-checked against the independent enumeration
\texttt{enum\_lf\_lex\_exact}.

\subsection{Optimal carbon tax and the cost--emission trade-off}
\label{sec:exp:tradeoff}

\Cref{tab:results_main} reports the eight validated base runs.

\begin{table}[ht]
  \centering
  \caption{Validated lexicographic \DBP runs on the base instances.}
  \label{tab:results_main}
  % !TeX root = ../main.tex
% ============================================================
%  tables/tab_results_main.tex
%  Source: results/analysis/tables/master_runs.csv (category=base)
% ============================================================

\renewcommand{\arraystretch}{1.15}
\setlength{\tabcolsep}{4pt}
\small
\begin{tabular}{@{}lrrrrrr@{}}
  \toprule
  \textbf{Instance}
    & $\tau^{\star}$
    & $E^{\star}$
    & $C^{\star}$
    & \textbf{Iter.}
    & \textbf{Rel. gap}
    & \textbf{Time (s)} \\
  \midrule
  STD\_I17\_M4\_T3 & 0.04 &  32{,}231 &   632 & 2 & 8.68e-09 &    36.7 \\
  STD\_I17\_M4\_T6 & 0.03 &  99{,}044 & 1{,}544 & 2 & 7.56e-10 &    51.8 \\
  STD\_I17\_M8\_T3 & 0.20 & 119{,}257 &   635 & 3 & 1.05e-06 &   127.0 \\
  STD\_I17\_M8\_T6 & 0.20 & 224{,}362 & 2{,}237 & 2 & 8.31e-10 & 1{,}936.0 \\
  STD\_I25\_M4\_T3 & 0.15 &  15{,}898 &   431 & 3 & 8.85e-12 &    37.6 \\
  STD\_I25\_M4\_T6 & 0.02 &  31{,}533 &   910 & 2 & 1.61e-09 &   138.0 \\
  STD\_I25\_M8\_T3 & 0.02 & 130{,}002 &   982 & 2 & 8.02e-16 &   152.6 \\
  STD\_I25\_M8\_T6 & 0.07 & 144{,}754 & 2{,}118 & 2 & 7.17e-10 & 2{,}367.0 \\
  \bottomrule
  \multicolumn{7}{@{}p{0.95\linewidth}@{}}{\footnotesize
    Notes: the eight solved base instances, method
    \texttt{lf\_lex\_exact}, under
    \path{results/runs/base/*/lf_lex_exact}. All rows have
    \texttt{success=True}, are \texttt{validated}, and terminate on the
    relative-gap criterion with gap $\le10^{-4}$. $\tau^{\star}$ is in
    USD/kg\,CO\textsubscript{2}e; $E^{\star}$ in 10\textsuperscript{3}
    kg\,CO\textsubscript{2}e and $C^{\star}$ in 10\textsuperscript{3}
    monetary units. The two M8\,T3/T6 rows with $\tau^{\star}=0.20$ sit
    at the base tax ceiling $\taxub=20$ and are re-examined under extended
    ranges in \Cref{sec:exp:tradeoff}.}\\
\end{tabular}

\end{table}

The regulator's optimal tax is strongly \emph{instance-specific}: across the
eight networks $\tau^{\star}$ spans the full base range, from $20$ to
$200$ USD/t\,CO\textsubscript{2}e. There is no single ``right'' carbon
price; the rate that minimises emissions for one network configuration can
be an order of magnitude away from the rate for another. Because the
follower carries binary facility, technology, capacity, and transport
decisions, its response shifts in discrete blocks rather than smoothly, so
nearby tax levels can induce the same network and identical emissions.

This blockiness has a sharp consequence at the high end. Two instances
(\texttt{STD\_I17\_M8\_T3} and \texttt{STD\_I17\_M8\_T6}) place
$\tau^{\star}$ exactly at the base ceiling of $200$ USD/t, signalling that the
base range may truncate the true optimum. The extended-range runs
(\Cref{app:tax_extended}) reveal a \emph{saturating} emission--tax response:
beyond a moderate rate, additional tax buys negligible further abatement. For \texttt{STD\_I25\_M8\_T3}, lifting
the ceiling moves the lexicographic optimum from $20$ to $360$ USD/t
while cutting emissions by only $0.04\%$ --- an eighteen-fold higher tax for a
rounding-error reduction. For \texttt{STD\_I25\_M4\_T3} the optimum moves from
$150$ to $310$ USD/t for a $0.016\%$ reduction, and for
\texttt{STD\_I17\_M8\_T3} emissions are flat to within solver tolerance across
$\tau\in[140,310]$ USD/t, so the emissions-minimising tax is not point-identified
at all but spread over a wide band.

This is the central caution of the study. A regulator whose \emph{sole}
objective is emissions minimisation is ill-posed on a saturating frontier: it
can rationalise an arbitrarily heavy tax for vanishing environmental gain.
The lexicographic structure of \eqref{eq:leader} --- minimise emissions, then
levy the lowest rate that attains the minimum --- is what disciplines this
pathology, selecting the bottom of the near-flat band; but it does so only up
to the solver's relative-gap tolerance. The policy reading is that the
emissions target should be paired with an explicit cost criterion (or a
minimum meaningful-abatement threshold) rather than pursued in isolation.

\subsection{What drives cost and emissions}
\label{sec:exp:sensitivity}

\Cref{tab:sensitivity} ranks the input factors by their normalized influence
on cost and on emissions.

\begin{table}[ht]
  \centering
  \caption{Normalized local sensitivity (OAT elasticities).}
  \label{tab:sensitivity}
  % !TeX root = ../main.tex
% ============================================================
%  tables/tab_sensitivity.tex
%  Source: results/analysis/tables/sensitivity_matrix.csv
%          (report SENSITIVITY_REPORT.md, section 1)
% ============================================================

\renewcommand{\arraystretch}{1.15}
\setlength{\tabcolsep}{6pt}
\small
\begin{tabular}{@{}lrr@{}}
  \toprule
  \textbf{Input factor}
    & \textbf{Cost elast.}
    & \textbf{Emissions elast.} \\
  \midrule
  Production capacity            & 0.58 & 0.04 \\
  Demand                         & 0.57 & 0.98 \\
  Production emission intensity  & 0.12 & 0.56 \\
  Return availability            & 0.09 & $<$0.01 \\
  Return processing cost         & 0.05 & $<$0.01 \\
  Scope-3 (supply) emissions     & 0.04 & 0.42 \\
  Bill-of-materials ratio        & 0.03 & $<$0.01 \\
  \midrule
  \makecell[l]{Warehouse cost, truck capacity,\\tax-range ceiling, air/sea cost}
                                 & $\le$0.02 & $<$0.01 \\
  \bottomrule
  \multicolumn{3}{@{}p{0.86\linewidth}@{}}{\footnotesize
    Notes: mean absolute normalized elasticities across the four screened
    base instances (I17/I25 $\times$ M4/M8, $T=3$); a value of $1.0$ means a
    $1\%$ input change moves the output by $1\%$. Structural perturbations
    without a scalar input change carry no elasticity; 31 of the screened
    perturbations fall below the $0.05\%$ solver-noise floor.}\\
\end{tabular}

\end{table}

The two outputs are governed by different factors. \emph{Cost} responds most
to production capacity (elasticity $0.58$) and demand ($0.57$); \emph{emissions}
respond almost one-for-one to demand ($0.98$) and strongly to the production
and scope-3 (upstream supply) emission intensities ($0.56$ and $0.42$). The
ranking is stable across network sizes --- demand is the top cost factor in
three of the four screened instances (production emission intensity in the
fourth) --- which strengthens the generality of the finding. The Morris screen
agrees that demand dominates, though as a single-trajectory screen on one
instance it is preliminary and we treat it as confirmatory only.

The economically important message is what does \emph{not} matter:
transport cost, truck capacity, warehouse cost, and the air/sea cost all carry
emissions elasticities below $0.01$. Emissions in these networks are
determined by \emph{what is produced and procured}, not by how goods are
moved. This directly motivates the design-lever experiment that follows.

\subsection{Design levers: technology versus transport}
\label{sec:exp:scenarios}

\Cref{tab:scenarios} contrasts the three design-lever scenarios against the
base instances.

\begin{table}[ht]
  \centering
  \caption{Optimal-response emissions under the design-lever scenarios.}
  \label{tab:scenarios}
  % !TeX root = ../main.tex
% ============================================================
%  tables/tab_scenarios.tex
%  Source: results/analysis/tables/scenario_2x2.csv
% ============================================================

\renewcommand{\arraystretch}{1.15}
\setlength{\tabcolsep}{6pt}
\small
\begin{tabular}{@{}lrrrr@{}}
  \toprule
  \textbf{Instance}
    & \textbf{Base $E^{\star}$}
    & \textbf{Technology}
    & \textbf{Transport}
    & \textbf{Combined} \\
  \midrule
  STD\_I17\_M4\_T3 &  32{,}231 & 22{,}083\,(\textminus31.5\%) &  32{,}242\,(+0.03\%) & 22{,}092\,(\textminus31.5\%) \\
  STD\_I17\_M8\_T3 & 119{,}257 & 80{,}440\,(\textminus32.5\%) & 119{,}298\,(+0.03\%) & 80{,}515\,(\textminus32.5\%) \\
  STD\_I25\_M4\_T3 &  15{,}898 & 11{,}049\,(\textminus30.5\%) &  15{,}920\,(+0.14\%) & 11{,}054\,(\textminus30.5\%) \\
  STD\_I25\_M8\_T3 & 130{,}002 & 96{,}661\,(\textminus25.6\%) & 130{,}023\,(+0.02\%) & 96{,}725\,(\textminus25.6\%) \\
  \bottomrule
  \multicolumn{5}{@{}p{0.95\linewidth}@{}}{\footnotesize
    Notes: optimal-response emissions $E^{\star}$ (10\textsuperscript{3}
    kg\,CO\textsubscript{2}e) under the base instance and three design-lever
    scenarios; parentheses give the change relative to the base. \emph{Technology}
    expands the production/remanufacturing technology menu, \emph{Transport}
    expands the transport-mode catalogue, and \emph{Combined} applies both.
    Expanding technology drives essentially all of the abatement; expanding
    transport modes alone is inert.}\\
\end{tabular}

\end{table}

Expanding the production and remanufacturing technology menu reduces optimal
emissions by $26\%$ to $33\%$ across all four instances. Expanding the
transport-mode catalogue alone changes emissions by essentially zero (at most
$+0.1\%$), and the combined scenario is indistinguishable from technology
alone. Cleaner production and remanufacturing capability, not modal choice, is
the binding abatement lever in these networks --- exactly the pattern the
elasticities predicted. The cost side is more nuanced: the technology menu can
raise cost slightly on the smaller I17 instances ($+2$ to $+4\%$) but sharply
\emph{reduces} both cost and emissions on the larger I25 networks (\eg
\texttt{STD\_I25\_M4\_T3}, where the combined scenario cuts cost by $42\%$
while abating $30.5\%$ of emissions), indicating that for larger networks the
expanded technology set unlocks configurations that are jointly cheaper and
cleaner.

\subsection{Solver behaviour and validation}
\label{sec:exp:scalability}

The findings above rest on the \DBP returning trustworthy solutions across the
entire experimental space. It does: all 95 committed runs are validated, the
largest relative gap is $1.9\times10^{-5}$ (and $1.0\times10^{-6}$ on the base
instances), none exceeds the target tolerance, and the procedure converges in
two to three iterations on average. \Cref{tab:comparison} profiles runtime by
planning horizon.

\begin{table}[ht]
  \centering
  \caption{Runtime profile of the validated base runs.}
  \label{tab:comparison}
  % !TeX root = ../main.tex
% ============================================================
%  tables/tab_comparison.tex
%  Source: results/analysis/tables/master_runs.csv (category=base)
% ============================================================

\renewcommand{\arraystretch}{1.2}
\begin{tabular}{@{}lrrrr@{}}
  \toprule
  \textbf{Planning horizon}
    & \textbf{Solved runs}
    & \textbf{Min (s)}
    & \textbf{Median (s)}
    & \textbf{Max (s)} \\
  \midrule
  $T=3$ & 4 & 36.7 &   82.3 &   152.6 \\
  $T=6$ & 4 & 51.8 & 1{,}037.0 & 2{,}367.0 \\
  \bottomrule
  \multicolumn{5}{@{}p{0.82\linewidth}@{}}{\footnotesize
    Notes: wall-clock time of the eight validated base runs in
    \Cref{tab:results_main}, grouped by planning horizon. The four $T=12$
    base instances are intractable within the compute budget and are not
    solved. Runtime rises steeply with the horizon and with the
    larger I25--M8 network.}\\
\end{tabular}

\end{table}

Runtime is driven by the planning horizon and the larger I25--M8 network: T3
instances solve in under three minutes, T6 instances in up to forty minutes,
and the most demanding run in the whole study --- the \texttt{STD\_I25\_M8\_T3}
combined scenario --- takes $1.86\times10^{4}$\,s ($\approx5.2$\,h). The four
T12 base instances do not solve within the compute budget, which bounds the
present study to the T3--T6 range and marks horizon length as the primary
computational bottleneck.
